\textbf{Phần 1. Trắc nghiệm lựa chọn}
\setcounterpageref{ex}{0}

\begin{ex} %[2D8V37]
Tổng $5 + 5 + 5 + 5 + 5 + 5$ được viết thành phép nhân nào dưới đây?
\choice
{$5 \times 5$}
{$6 \times 5$}
{\True $5 \times 6$}
{$5 \times 7$}
\loigiai{
Số 5 được lấy 6 lần nên viết là $5 \times 6$.
\textbf{Đáp án của câu hỏi là C}
}
\end{ex}

\begin{ex} %[2D8V37]
Một bàn học có $2$ bạn học sinh ngồi. Hỏi $8$ bàn học như thế có tất cả bao nhiêu bạn học sinh ngồi?
\choice
{\True $16$ bạn}
{$14$ bạn}
{$18$ bạn}
{$10$ bạn}
\loigiai{
$2 \times 8 = 16$ (bạn).
\textbf{Đáp án của câu hỏi là A}
}
\end{ex}

\begin{ex} %[2D8V37]
Phép nhân $4 \times 3$ biểu diễn tổng của
\choice
{$4$ số hạng $3$}
{$4$ cộng với $3$}
{$4$ trừ đi $3$}
{\True $3$ số hạng $4$ cộng lại với nhau}
\loigiai{
$4 \times 3 = 4 + 4 + 4$, tức là 3 số hạng 4 cộng lại.
\textbf{Đáp án của câu hỏi là D}
}
\end{ex}

\begin{ex} %[2D8V37]
Mỗi con thỏ có $4$ cái chân. Hỏi $5$ con thỏ có tất cả bao nhiêu cái chân?
\choice
{$16$ cái chân}
{\True $20$ cái chân}
{$24$ cái chân}
{$15$ cái chân}
\loigiai{
$4 \times 5 = 20$ (cái chân).
\textbf{Đáp án của câu hỏi là B}
}
\end{ex}

\begin{ex} %[2D8V37]
Tính giá trị của biểu thức: $3 \times 4 + 15$.
\choice
{$25$}
{$26$}
{\True $27$}
{$28$}
\loigiai{
Thực hiện nhân trước: $3 \times 4 = 12$; $12 + 15 = 27$.
\textbf{Đáp án của câu hỏi là C}
}
\end{ex}

\begin{ex} %[2D8V37]
Mỗi gói kẹo có $5$ cái kẹo. Bạn Lan mua $4$ gói kẹo và được tặng thêm $2$ cái kẹo nữa. Lan có tất cả bao nhiêu cái kẹo?
\choice
{\True $22$ cái kẹo}
{$20$ cái kẹo}
{$24$ cái kẹo}
{$18$ cái kẹo}
\loigiai{
$5 \times 4 + 2 = 20 + 2 = 22$ (cái kẹo).
\textbf{Đáp án của câu hỏi là A}
}
\end{ex}

\begin{ex} %[2D8V37]
Điền dấu thích hợp: $2 \times 7 \dots 3 \times 5$.
\choice
{$>$}
{$=$}
{Không so sánh được}
{\True $<$}
\loigiai{
$2 \times 7 = 14$; $3 \times 5 = 15$. Vì $14 < 15$ nên điền dấu $<$. 
\textbf{Đáp án của câu hỏi là D}
}
\end{ex}

\begin{ex} %[2D8V37]
Một xe đạp có $2$ bánh xe. Hỏi $9$ xe đạp như thế có bao nhiêu bánh xe?
\choice
{$16$ bánh xe}
{\True $18$ bánh xe}
{$20$ bánh xe}
{$14$ bánh xe}
\loigiai{
$2 \times 9 = 18$ (bánh xe).
\textbf{Đáp án của câu hỏi là B}
}
\end{ex}

\begin{ex} %[2D8V37]
Viết tổng $2 + 2 + 2 + 2 + 2 + 2 + 2$ dưới dạng tích:
\choice
{$2 \times 6$}
{$7 \times 2$}
{\True $2 \times 7$}
{$2 \times 8$}
\loigiai{
Số 2 được lấy 7 lần: $2 \times 7$.
\textbf{Đáp án của câu hỏi là C}
}
\end{ex}

\begin{ex} %[2D8V37]
Tìm một số biết rằng lấy số đó trừ đi $12$ thì bằng tích của $2$ và $6$.
\choice
{\True $24$}
{$22$}
{$26$}
{$20$}
\loigiai{
Tích của 2 và 6 là $2 \times 6 = 12$. Số đó là: $12 + 12 = 24$.
\textbf{Đáp án của câu hỏi là A}
}
\end{ex}

\begin{ex} %[2D8V37]
Một tuần lễ có $7$ ngày. Hỏi $2$ tuần lễ có bao nhiêu ngày?
\choice
{$12$ ngày}
{$16$ ngày}
{$10$ ngày}
{\True $14$ ngày}
\loigiai{
$7 \times 2 = 14$ (ngày).
\textbf{Đáp án của câu hỏi là D}
}
\end{ex}

\begin{ex} %[2D8V37]
Có $4$ hàng cây, mỗi hàng có $5$ cây. Hỏi có tất cả bao nhiêu cây?
\choice
{$18$ cây}
{\True $20$ cây}
{$24$ cây}
{$16$ cây}
\loigiai{
$5 \times 4 = 20$ (cây).
\textbf{Đáp án của câu hỏi là B}
}
\end{ex}

\textbf{Phần 2. Trắc nghiệm đúng sai}
\setcounterpageref{ex}{0}

\begin{ex} %[2D8V37]
Xét tính đúng sai của các nhận định về phép nhân:
\choiceTF
{\True Phép nhân là phép cộng các số hạng bằng nhau.}
{Biểu thức $3 + 3 + 3 + 4$ có thể viết thành $3 \times 4$.}
{\True Kết quả của $2 \times 5$ bằng kết quả của $5 \times 2$.}
{\True Tích của bất kỳ số nào với 1 cũng bằng chính số đó.}
\loigiai{
\begin{itemize}
    \item Ý a) Đúng: Đúng theo bản chất phép nhân.
    \item Ý b) Sai: Các số hạng không bằng nhau (có số 4) nên không thể viết thành $3 \times 4$.
    \item Ý c) Đúng: Tính chất giao hoán của phép nhân: $2 \times 5 = 5 \times 2 = 10$.
    \item Ý d) Đúng: Đúng.
\end{itemize}
\begin{itemize}
    \item a) \textbf{Đúng}.
    \item b) \textbf{Sai}.
    \item c) \textbf{Đúng}.
    \item d) \textbf{Đúng}.
\end{itemize}
}
\end{ex}

\begin{ex} %[2D8V37]
Có 4 đĩa cam, mỗi đĩa có 3 quả cam:
\choiceTF
{\True Tổng số cam là 12 quả.}
{Phép nhân biểu diễn số cam trên là $4 \times 3$.}
{\True Nếu ăn đi 2 quả cam thì còn lại 10 quả cam.}
{Nếu thêm 1 đĩa cam nữa thì có tất cả 16 quả cam.}
\loigiai{
\begin{itemize}
    \item Ý a) Đúng: $3 \times 4 = 12$ quả.
    \item Ý b) Sai: Số quả cam ở mỗi đĩa là 3, được lấy 4 lần nên viết chuẩn là $3 \times 4$.
    \item Ý c) Đúng: $12 - 2 = 10$ quả.
    \item Ý d) Sai: Thêm 1 đĩa có thêm 3 quả: $12 + 3 = 15$ quả (không phải 16 quả).
\end{itemize}
\begin{itemize}
    \item a) \textbf{Đúng}.
    \item b) \textbf{Sai}.
    \item c) \textbf{Đúng}.
    \item d) \textbf{Sai}.
\end{itemize}
}
\end{ex}

\textbf{Phần 3. Trả lời ngắn}
\setcounterpageref{ex}{0}

\begin{ex} %[2D8V37]
Tính giá trị của biểu thức: $2 \times 6 + 18$.
\shortans{30}
\loigiai{
\textbf{Đáp số:} $30$\\
$2 \times 6 = 12$; $12 + 18 = 30$.
}
\end{ex}

\begin{ex} %[2D8V37]
Mỗi con bò có 4 cái chân. Hỏi 6 con bò có tất cả bao nhiêu cái chân?
\shortans{24}
\loigiai{
\textbf{Đáp số:} $24$\\
$4 \times 6 = 24$ cái chân.
}
\end{ex}

\begin{ex} %[2D8V37]
Có 5 bạn học sinh, mỗi bạn có 3 quyển vở. Hỏi cả 5 bạn có tất cả bao nhiêu quyển vở?
\shortans{15}
\loigiai{
\textbf{Đáp số:} $15$\\
$3 \times 5 = 15$ quyển vở.
}
\end{ex}

\begin{ex} %[2D8V37]
Tính nhanh tổng sau bằng phép nhân: $5 + 5 + 5 + 5 + 5$.
\shortans{25}
\loigiai{
\textbf{Đáp số:} $25$\\
$5 \times 5 = 25$.
}
\end{ex}

\begin{ex} %[2D8V37]
Một xe ô tô con có 4 bánh xe. Hỏi 8 xe ô tô con như thế có bao nhiêu bánh xe?
\shortans{32}
\loigiai{
\textbf{Đáp số:} $32$\\
$4 \times 8 = 32$ bánh xe.
}
\end{ex}

\begin{ex} %[2D8V37]
Tìm một số biết rằng lấy số đó chia cho 2 thì được kết quả bằng 7.
\shortans{14}
\loigiai{
\textbf{Đáp số:} $14$\\
$7 \times 2 = 14$.
}
\end{ex}
